\chapter{Kohomologia snopów z nośnikami i podstawy dualności Verdiera}
\label{ch:snopy-verdier-v12}

\section{Dlaczego teoria snopów jest potrzebna?}
W rozdziale~\ref{ch:compact-bm-v11} spotkaliśmy dwie postacie dualności
Poincarégo dla niezwartych rozmaitości: jedną z kohomologią o nośnikach
zwartych, drugą z homologią Borel--Moore'a. Ich elementarne dowody
wykorzystywały lokalnie skończone triangulacje. Tymczasem w geometrii
algebraicznej często pracujemy na przestrzeniach osobliwych, a zamiast
jednej grupy współczynników pojawiają się dane zmieniające się lokalnie.
Potrzebujemy konstrukcji, która jednocześnie pamięta (i) współczynniki
lokalne, (ii) warunki na nośniki oraz (iii) sposób, w jaki lokalne
dane sklejają się w globalne.

Snop jest właśnie uporządkowanym systemem takich danych. Kohomologia
snopów mierzy przeszkody w ich globalnym sklejaniu, zaś dualność Verdiera
jest snopowym uogólnieniem dualności Poincarégo. Rozpoczniemy od
konstrukcji elementarnych, a twierdzenia wymagające kategorii pochodnych
wyraźnie oddzielimy od wyników, których dowody przedstawimy tutaj.

W całym rozdziale $k$ jest ustalonym ciałem, chyba że jawnie podamy inny
pierścień współczynników. Przestrzenie są Hausdorffa i lokalnie zwarte,
a w twierdzeniach porównawczych i o nośnikach dodatkowo parazwarte.
Kohomologia jest indeksowana \emph{kohomologicznie}: różniczka podnosi
stopień o jeden. Zapis $\mathcal F$ oznacza snop $k$-przestrzeni wektorowych.

\section{Przedsnopy, snopy i włókna snopu}
\begin{definicja}[Przedsnop i warunek snopowy]
\label{def:sheaf-v12}
Niech $X$ będzie przestrzenią topologiczną. Przedsnopem
$k$-przestrzeni wektorowych $\mathcal F$ na $X$ nazywamy przyporządkowanie
każdemu otwartemu $U\subseteq X$ przestrzeni $\mathcal F(U)$ i każdemu
zawieraniu $V\subseteq U$ liniowej mapy ograniczenia
$\rho^U_V:\mathcal F(U)\to\mathcal F(V)$, przy czym
$\rho^U_U=\id$ i $\rho^V_W\rho^U_V=\rho^U_W$ dla $W\subseteq V\subseteq U$.
Przedsnop jest \emph{snopem}, jeśli dla każdego otwartego $U$ oraz
pokrycia $U=\bigcup_iU_i$ zachodzą:
\begin{enumerate}[label=(\roman*)]
\item jeśli $s,t\in\mathcal F(U)$ mają równe ograniczenia do każdego $U_i$,
  to $s=t$ (jednoznaczność);
\item jeśli $s_i\in\mathcal F(U_i)$ i
  $s_i|_{U_i\cap U_j}=s_j|_{U_i\cap U_j}$ dla wszystkich $i,j$,
  istnieje $s\in\mathcal F(U)$ o $s|_{U_i}=s_i$ (istnienie sklejenia).
\end{enumerate}
\end{definicja}

\begin{przyklad}[Snopy funkcji, form i rozwiązań lokalnych]
Na gładkiej rozmaitości $M$ przyporządkowania
$U\mapsto C^\infty(U)$ oraz $U\mapsto\Omega^q(U)$ są snopami,
podobnie jak $U\mapsto\mathcal O(U)$ dla holomorficznych funkcji na
zespolonej rozmaitości. Dla ustalonego operatora różniczkowego $P$
przestrzenie $\{f\in C^\infty(U):Pf=0\}$ również tworzą snop
(jeśli operator jest zgodny z ograniczeniami).
Natomiast przyporządkowanie $U\mapsto H^q(U;k)$ dla $q>0$ nie jest
na ogół snopem: globalna klasa kohomologii może znikać na wszystkich
zbiorach drobnego pokrycia, nie będąc zerowa globalnie.
Na przykład generator $H^1(S^1;k)$ znika po ograniczeniu
do każdej z dwóch dostatecznie dużych otwartych łuków.
\end{przyklad}

\begin{definicja}[Włókno snopu w punkcie]
Dla snopu $\mathcal F$ i $x\in X$ jego \emph{włóknem} (stalkiem)
nazywamy granicę prostą
$\mathcal F_x=\varinjlim_{x\in U}\mathcal F(U)$, gdzie przejścia
są ograniczeniami, a $U$ przebiegają otwarte otoczenia $x$.
Element $\mathcal F_x$ jest klasą pary $(U,s)$, $x\in U$,
z utożsamieniem dwóch par, jeśli sekcje zgadzają się
na pewnym mniejszym otoczeniu $x$.
\end{definicja}

\begin{lemat}[Dokładność sprawdza się na włóknach]
\label{lem:stalk-exact-v12}
Ciąg morfizmów snopów $k$-przestrzeni wektorowych
$\mathcal A\xrightarrow{u}\mathcal B\xrightarrow{v}\mathcal C$
jest dokładny (czyli $\operatorname{im}u=\ker v$ jako snopy)
wtedy i tylko wtedy, gdy ciąg
$\mathcal A_x\xrightarrow{u_x}\mathcal B_x\xrightarrow{v_x}\mathcal C_x$
jest dokładny dla każdego $x\in X$.
\end{lemat}
\begin{proof}
Warunek jądra jest lokalny, ponieważ dla sekcji $b$ równość $v(b)=0$
można sprawdzać na wszystkich jej kiełkach. Obraz \emph{snopowy}
morfizmu $u$ składa się z sekcji, które lokalnie należą do obrazu $u$.
Zatem kiełek $b_x$ należy do $\operatorname{im}u_x$ dokładnie wtedy,
gdy po zmniejszeniu otoczenia $x$ sekcja $b$ ma lokalny przeciwobraz.
To równoważność dokładności w punkcie i dokładności snopowej.
Nie wynika stąd, że ciąg przestrzeni sekcji globalnych jest dokładny:
lokalne przeciwobrazy mogą nie sklejać się globalnie.
\end{proof}

\begin{przyklad}[Snop stały a system lokalny]
Na spójnej przestrzeni lokalnie łukowo spójnej snop $k_X$
przyporządkowuje otwartemu $U$ \emph{lokalnie stałe} funkcje $U\to k$;
nie należy go mylić z przedsnopem $U\mapsto k$ nawet wtedy,
gdy $U$ jest niespójne. Każdy system lokalny o włóknie $k^r$
można otrzymać z reprezentacji
$\pi_1(X,x_0)\to\operatorname{GL}_r(k)$,
jeśli $X$ jest spójna, lokalnie łukowo spójna i półlokalnie jednospójna.
\end{przyklad}

\section{Nośnik sekcji i przedłużenie przez zero}
\begin{definicja}[Sekcja o nośniku zwartym]
\label{def:gamma-c-v12}
Niech $X$ będzie lokalnie zwarta i Hausdorffa, a $\mathcal F$
niech będzie snopem na $X$. Dla $s\in\mathcal F(X)$ określmy
$\operatorname{supp}(s)=\{x\in X:s_x\ne0\}$
(zbiór ten jest już domknięty na mocy warunku snopowego).
Przestrzeń sekcji o nośniku zwartym oznaczamy
\[\Gamma_c(X,\mathcal F)
  :=\{s\in\mathcal F(X):\operatorname{supp}(s)\text{ jest zwarty}\}.\]
Dla domkniętego $Z\subseteq X$ analogicznie oznaczamy
$\Gamma_Z(X,\mathcal F)=\{s:\operatorname{supp}(s)\subseteq Z\}$.
\end{definicja}

\begin{uwaga}[Skąd domkniętość nośnika?]
Zbiór punktów $x$ o $s_x=0$ jest otwarty: zerowy kiełek oznacza,
że na pewnym otoczeniu $x$ sekcja $s$ jest zerowa. Jego dopełnienie
jest więc domknięte. Dla snopu funkcji jest to zwykłe domknięcie
zbioru, na którym funkcja jest różna od zera.
\end{uwaga}

\begin{definicja}[Przedłużenie przez zero]
\label{def:jbang-v12}
Niech $j:U\hookrightarrow X$ będzie zanurzeniem \emph{otwartym}
i $\mathcal F$ snopem na $U$. Przez $j_!\mathcal F$ oznaczamy snop,
który ma na otwartym $V\subseteq X$ następujące sekcje:
\[(j_!\mathcal F)(V)=\{s\in\mathcal F(V\cap U):
\operatorname{supp}_{V\cap U}(s)\text{ jest domknięty w }V\}.\]
W szczególności sekcje te dają się przedłużyć przez zero
poza $U$. Na włóknach:
\[
(j_!\mathcal F)_x\cong
\begin{cases}\mathcal F_x,&x\in U,\\0,&x\notin U.\end{cases}
\]
Natomiast obraz bezpośredni $j_*\mathcal F$ definiuje wzór
$(j_*\mathcal F)(V)=\mathcal F(V\cap U)$ dla otwartych $V\subseteq X$.
\end{definicja}

\begin{przyklad}[Dlaczego $j_!$ różni się od $j_*$?]
Niech $X=\R$, $U=(0,1)$ i $\mathcal F=k_U$. Dla $x=0$ mamy
$(j_!k_U)_0=0$, ale $(j_*k_U)_0\cong k$, gdyż dla małego
$\varepsilon>0$ funkcje lokalnie stałe na $(0,\varepsilon)$
mają jedną swobodną wartość. Istotne jest więc, czy
informację dopuszczamy \emph{aż do brzegu}, czy znikającą przy nim.
\end{przyklad}

\begin{center}
\begin{tikzpicture}[x=1cm,y=1cm,font=\small]
\draw[-{Stealth[length=2mm]},thick] (-3.9,0)--(4.0,0);
\node[below] at (-3.4,0) {$X=\R$};
\draw[line width=2.3pt,manifoldblue] (-1.9,0)--(2.0,0);
\draw[fill=white,thick] (-1.9,0) circle (2pt);
\draw[fill=white,thick] (2.0,0) circle (2pt);
\node[above,manifoldblue] at (0,0.12) {$U=(0,1)$ (schematycznie)};
\node[below] at (-1.9,-0.14) {$0$};
\node[below] at (2.0,-0.14) {$1$};
\node[below,align=center] at (0,-0.70) {Na brzegu: $(j_!k_U)_x=0$,\\ ale $(j_*k_U)_x\cong k$.};
\end{tikzpicture}
\par\smallskip\small Przedłużenie przez zero zachowuje dane wewnątrz $U$,
ale usuwa ich kiełki na brzegu otwartego podzbioru.
\end{center}

\begin{stwierdzenie}[Sekcje na otwartym podzbiorze]
\label{prop:jbang-sections-v12}
Przy założeniach definicji~\ref{def:jbang-v12}, jeśli $X$ jest lokalnie
zwarta i Hausdorffa, to zachodzi naturalna identyfikacja
$\Gamma_c(X,j_!\mathcal F)\cong\Gamma_c(U,\mathcal F)$.
\end{stwierdzenie}
\begin{proof}
Sekcja $j_!\mathcal F$ jest na $X\setminus U$ lokalnie zerowa,
a więc odpowiada sekcji $\mathcal F$ na $U$ z zerowym przedłużeniem.
Jeśli jej nośnik jest zwarty w $X$, to leży w $U$ i jest
zwarty również w $U$; w drugą stronę nośnik zwarty w $U$
jest zwarty w $X$ (z ciągłości włożenia), a przedłużenie
przez zero określa sekcję $j_!\mathcal F$.
Konstrukcje są odwrotne i naturalne.
\end{proof}

\section{Kohomologia snopów i kohomologia z nośnikami}
\begin{definicja}[Funktory pochodne sekcji]
\label{def:sheaf-cohom-v12}
Dla snopu $\mathcal F$ na $X$ wybieramy rezolwentę iniektywną
$0\to\mathcal F\to\mathcal I^0\to\mathcal I^1\to\cdots$
w kategorii snopów $k$-przestrzeni wektorowych. Definiujemy
\[
H^q(X,\mathcal F)=H^q\bigl(\Gamma(X,\mathcal I^\bullet)\bigr),
\quad H_c^q(X,\mathcal F)
=H^q\bigl(\Gamma_c(X,\mathcal I^\bullet)\bigr),
\quad H_Z^q(X,\mathcal F)
=H^q\bigl(\Gamma_Z(X,\mathcal I^\bullet)\bigr).
\]
Są to odpowiednio kohomologia snopu, kohomologia z nośnikami
zwartymi oraz kohomologia lokalna z nośnikiem w $Z$.
Ich niezależność od wyboru rezolwenty wynika z twierdzenia
porównawczego dla rezolwent iniektywnych.
\end{definicja}

\begin{uwaga}[Co zakładamy z algebry homologicznej?]
Istnienie dostatecznej liczby snopów iniektywnych i twierdzenie
o niezależności rezolwenty są standardowymi rezultatami
algebry homologicznej w kategorii snopów. Konstrukcja powyższa
nie \emph{dowodzi} ich istnienia; szczegóły kategorii
pochodnej rozwiniemy przed dualnością Verdiera.
\end{uwaga}

\begin{lemat}[Ciąg Mayera--Vietorisa ze zwartymi nośnikami]
\label{lem:sheaf-mv-v12}
Niech $X$ będzie lokalnie zwartą przestrzenią Hausdorffa,
$X=U\cup V$ pokryciem otwartym, a $\mathcal F$ snopem
$k$-przestrzeni wektorowych na $X$. Zachodzi naturalny
długi ciąg dokładny
\[
\cdots\longrightarrow H_c^q(U\cap V,\mathcal F)
\longrightarrow H_c^q(U,\mathcal F)\oplus H_c^q(V,\mathcal F)
\longrightarrow H_c^q(X,\mathcal F)
\longrightarrow H_c^{q+1}(U\cap V,\mathcal F)\longrightarrow\cdots,
\]
gdzie wszędzie po lewej stosujemy ograniczenie snopu do
wskazanego otwartego podzbioru.
\end{lemat}
\begin{proof}
Niech $a:U\cap V\hookrightarrow X$, $j:U\hookrightarrow X$
i $l:V\hookrightarrow X$ oznaczają zanurzenia otwarte.
Sprawdzenie na włóknach (Lemat~\ref{lem:stalk-exact-v12}) daje krótki
ciąg dokładny snopów
\[
0\longrightarrow a_!(\mathcal F|_{U\cap V})
\xrightarrow{\,s\mapsto(s,-s)\,}
 j_!(\mathcal F|_U)\oplus l_!(\mathcal F|_V)
\xrightarrow{\,(s,t)\mapsto s+t\,}\mathcal F\longrightarrow0.
\]
Na przecięciu jest to $0\to F\to F\oplus F\to F\to0$,
a poza przecięciem identyczność na odpowiednim składniku.
Ponieważ przedłużenie przez zero przy otwartym zanurzeniu
jest dokładne (co także wynika ze sprawdzenia na włóknach),
ciąg pozostaje dokładny po zastosowaniu rezolwent w kategorii
pochodnej. Następnie stosujemy funktor $R\Gamma_c(X,-)$.
Przez Stwierdzenie~\ref{prop:jbang-sections-v12} i dokładność
$j_!$ identyfikujemy otrzymane człony z
$R\Gamma_c(U\cap V,-)$, $R\Gamma_c(U,-)$ i $R\Gamma_c(V,-)$.
Długi ciąg kohomologii trójkąta dokładnego daje tezę.
\end{proof}

\begin{przyklad}[Obliczenie $H_c^*(\R;k)$]
Niech $U=(-\infty,1)$, $V=(0,\infty)$, więc $U\cap V=(0,1)$
i $\R=U\cup V$. Każdy z tych trzech zbiorów jest homeomorficzny
z $\R$, dlatego w tym przykładzie ciąg Mayera--Vietorisa
sam nie oblicza nieznanej grupy. Zamiast kołowego rozumowania
używamy definicji względnej z
Definicji~\ref{def:compact-support-cohom-v11}:
$H_c^q(\R;k)\cong H^q([-1,1],\{-1,1\};k)$
przez wycięcie i wyczerpanie przedziałami. Ciąg względny daje
$H_c^0(\R;k)=0$ i $H_c^1(\R;k)\cong k$.
Poniżej odzyskamy ten sam wynik z rezolwenty de Rhama.
\end{przyklad}

\section{Od kohomologii \v{C}echa do kohomologii snopów}
\label{sec:cech-sheaves-v12}
Nie należy mylić \emph{homologii \v{C}echa} z rozdziału o
nerwach pokryć z \emph{kohomologią \v{C}echa snopu}, której
współczynniki stanowią sekcje snopu na przecięciach otwartych.

\begin{definicja}[Kołańcuchy \v{C}echa ze współczynnikami w snopie]
Niech $\mathcal U=(U_i)_{i\in I}$ będzie uporządkowanym pokryciem
otwartym $X$, a $\mathcal F$ snopem abelowym. Definiujemy
\[
\check C^q(\mathcal U,\mathcal F)
=\prod_{i_0<\cdots<i_q}\mathcal F(U_{i_0}\cap\cdots\cap U_{i_q}),
\qquad
(\delta c)_{i_0\cdots i_{q+1}}
=\sum_{j=0}^{q+1}(-1)^j
c_{i_0\cdots\widehat{i_j}\cdots i_{q+1}}
\big|_{U_{i_0}\cap\cdots\cap U_{i_{q+1}}}.
\]
Niepuste przecięcia są przyjmowane w indeksie; dla pustego
przecięcia przestrzeń sekcji jest zerowa. Kohomologię tego
kompleksu oznaczamy $\check H^q(\mathcal U,\mathcal F)$.
\end{definicja}

\begin{stwierdzenie}[Pierwsza przeszkoda sklejania]
$\check H^0(\mathcal U,\mathcal F)=\mathcal F(X)$,
gdy $\mathcal F$ jest snopem i $\mathcal U$ pokrywa $X$.
\end{stwierdzenie}
\begin{proof}
Zero-kocykl to rodzina sekcji $(s_i)$ spełniająca
$s_j|_{U_i\cap U_j}=s_i|_{U_i\cap U_j}$.
Warunek snopowy daje dokładnie jedną sekcję na całym $X$.
W standardowym kompleksie \v{C}echa nie ma kołańcuchów stopnia $-1$,
zatem w stopniu zerowym nie bierzemy ilorazu przez żadne kobrzegi.
\end{proof}

\begin{twierdzenie}[Kiedy \v{C}ech oblicza kohomologię snopów?]
\label{thm:leray-sheaf-v12}
Jeżeli dla każdego niepustego skończonego przecięcia
$U_{i_0}\cap\cdots\cap U_{i_q}$ oraz każdego $r>0$ mamy
$H^r(U_{i_0}\cap\cdots\cap U_{i_q},\mathcal F)=0$,
to naturalna mapa
$\check H^m(\mathcal U,\mathcal F)\to H^m(X,\mathcal F)$
jest izomorfizmem dla każdego $m\ge0$.
\end{twierdzenie}
\begin{proof}[Podwójny kompleks i dwie filtracje]
Niech $\mathcal F\to\mathcal I^\bullet$ będzie rezolwentą
iniektywną. Dla uporządkowanego pokrycia $\mathcal U$
wprowadzamy dwukompleks
$C^{p,q}=\prod_{i_0<\cdots<i_p}
\Gamma(U_{i_0\cdots i_p},\mathcal I^q)$ dla $p,q\ge0$,
z różniczką całkowitą
$D=\delta+(-1)^p d_{\mathcal I}$ na $C^{p,q}$.
Obie filtracje są w pierwszej ćwiartce, zatem w każdym
stopniu całkowitym występuje skończenie wiele par $(p,q)$;
standardowe twierdzenie o zbieżności ma tu zastosowanie.

Najpierw liczymy kohomologię w kierunku $p$. Snop
iniektywny jest wiotki (flabby): ograniczenie jego sekcji
do zbioru otwartego jest surjektywne. Dla wiotkiego snopu
rozszerzony kompleks \v{C}echa danego pokrycia ma
zanik kohomologii w dodatnich stopniach. Istotne jest,
że chodzi tutaj o \emph{wiotkość} snopów $\mathcal I^q$,
a nie o sam warunek snopowy, który zapewnia tylko
dokładność w stopniu zerowym. Wobec tego pierwsza
filtracja daje kohomologię kompleksu
$\Gamma(X,\mathcal I^\bullet)$, czyli $H^*(X,\mathcal F)$.

W drugim kierunku, ponieważ iniektywne snopy są
acykliczne na dowolnym otwartym podzbiorze,
$H^q(\Gamma(U_{i_0\cdots i_p},\mathcal I^\bullet))
=H^q(U_{i_0\cdots i_p},\mathcal F)$.
Iloczyny przestrzeni wektorowych są dokładne,
więc $E_1^{p,q}=\prod_{i_0<\cdots<i_p}
H^q(U_{i_0\cdots i_p},\mathcal F)$.
Z hipotezy wynika zanik dla $q>0$.
Wiersz $q=0$ jest dokładnie kompleksem
$\check C^\bullet(\mathcal U,\mathcal F)$.
Dlatego drugi ciąg spektralny degeneruje po obliczeniu
kohomologii tego wiersza. Obie filtracje obliczają
kohomologię tego samego totalnego kompleksu;
indukowana mapa porównawcza jest zatem izomorfizmem.
Użyliśmy przy tym standardowego lematu o wiotkości
snopów iniektywnych, wyraźnie oddzielonego od rachunku
z dwiema filtracjami.
\end{proof}

\begin{przyklad}[Okrąg i jego lokalny system]
Niech $L$ będzie systemem lokalnym o włóknie $V=k^r$
na $S^1$, o monodromii $A\in\operatorname{GL}(V)$.
Jednokomórkowy model okręgu z odpowiednim skręceniem
współczynników ma kompleks kołańcuchów
$V\xrightarrow{A-I}V$ w stopniach $0$ i $1$.
Dlatego
$H^0(S^1;L)=\ker(A-I)$ oraz
$H^1(S^1;L)=\operatorname{coker}(A-I)$.
Na przykład dla $V=\R$ i $A=-1$ obie grupy są zerowe;
przy $A=1$ obie wynoszą $\R$.
To pokazuje, dlaczego znajomość samych włókien $L_x$
nie wystarcza: trzeba również znać ich transport.
\end{przyklad}

\section{Rezolwenta de Rhama i związki z poprzednim rozdziałem}
\begin{twierdzenie}[Twierdzenie de Rhama ze zwartymi nośnikami — dowód snopowy]
\label{thm:de-rham-sheaf-v12}
Niech $M$ będzie gładką drugoprzeliczalną rozmaitością,
bez brzegu. Wówczas naturalne całkowanie form wyznacza
$H^q(\Omega_c^\bullet(M),\dd)\cong H_c^q(M;\R)$,
gdzie prawa strona jest singularną kohomologią z nośnikami
zwartymi z Definicji~\ref{def:compact-support-cohom-v11}.
\end{twierdzenie}
\begin{proof}
Rozpatrzmy ciąg snopów
$0\to\R_M\to\Omega_M^0\xrightarrow{\dd}\Omega_M^1
\xrightarrow{\dd}\cdots$. Jest on dokładny na włóknach:
lokalny lemat Poincarégo dowodzi dokładności w dodatnich
stopniach, a w stopniu zero jądrem $\dd$ są funkcje lokalnie
stałe. Dla każdej formy i każdego otwartego pokrycia
istnieje gładki podział jedności podporządkowany
lokalnie skończonemu uszczegółowieniu; mnożenie form
przez jego funkcje daje operatory lokalizujące.
Dlatego snopy $\Omega_M^q$ są \emph{drobne} (fine).

Dla przestrzeni lokalnie zwartych i parazwartych snopy
drobne są acykliczne względem $\Gamma_c$. Wyjaśnijmy
mechanizm tego standardowego lematu: w kompleksie
\v{C}echa sekcji o nośnikach zwartych podział jedności
pozwala zbudować homotopię kontrahującą dodatnie
stopnie. Każda sekcja o nośniku zwartym spotyka tylko
skończenie wiele elementów lokalnie skończonego pokrycia,
więc operatory zachowują warunek zwartości nośnika;
na przestrzeni parazwartej porównanie rezolwenty
\v{C}echa z rezolwentą iniektywną daje zanik wyższych
funktorów pochodnych. W efekcie kompleks
$\Gamma_c(M,\Omega_M^\bullet)=\Omega_c^\bullet(M)$
oblicza $R\Gamma_c(M,\R_M)$.

Dla rozmaitości lokalnie kontraktowalnej porównanie
kohomologii stałego snopu z kohomologią singularną
(respektujące nośniki) utożsamia
$H_c^q(M,\R_M)$ z $H_c^q(M;\R)$.
Trzeba doprecyzować, co znaczy tutaj „całkowanie form”
w przypadku nośników zwartych. Dla ustalonego kompaktu
$K\subset M$ całkowanie po gładkich sympleksach daje
kołańcuch reprezentujący klasę względną
$H^q(M,M\setminus K;\R)$: forma o nośniku zawartym
w $K$ zeruje się na sympleksach o obrazie w $M\setminus K$.
Po przejściu do granicy prostej po kompaktach $K$
otrzymujemy mapę do $H_c^q(M;\R)$.
Jej zgodność z izomorfizmem snopowym wynika z lokalnego
lematu Poincarégo i naturalności konstrukcji względem
ograniczeń i ciągów Mayera--Vietorisa.
Stosujemy jako jawne twierdzenie porównawcze
identyfikację pochodnej kohomologii stałego snopu
z singularną na rozmaitościach; samo stwierdzenie,
że formy całkujemy po sympleksach, nie jest niezależnym
dowodem tego twierdzenia.
\end{proof}

\section{Kategorie pochodne: tylko niezbędny aparat}
\begin{definicja}[Kompleks i przesunięcie]
Dla kompleksu kohomologicznego $(A^\bullet,d)$ ustalamy
$A[n]^q=A^{q+n}$ oraz $d_{A[n]}=(-1)^n d_A$.
Morfizm kompleksów jest \emph{quasi-izomorfizmem}, gdy
indukuje izomorfizmy wszystkich grup kohomologii.
Kategoria pochodna $D(X)$ jest otrzymywana z kategorii
kompleksów snopów przez formalne odwrócenie wszystkich
quasi-izomorfizmów.
\end{definicja}

\begin{przyklad}[Odczytywanie przesunięcia]
Jeżeli $V$ jest przestrzenią wektorową skoncentrowaną
w stopniu $0$, to $V[n]$ ma jedyną niezerową kohomologię
w stopniu $-n$. W szczególności $\R_M[n]$ nie jest
zwykłym snopem ustawionym w stopniu $n$, tylko
kompleksem przesuniętym \emph{w lewo}.
\end{przyklad}

\begin{definicja}[Pochodne funktory Hom]
Przez $R\!\operatorname{Hom}_k(A^\bullet,B^\bullet)$
oznaczamy pochodny kompleks Hom, a przez
$R\mathcal Hom(\mathcal F^\bullet,\mathcal G^\bullet)$
pochodny snop wewnętrznych homomorfizmów.
Definiuje się je za pomocą stosownych rezolwent
iniektywnych; ich kohomologie mierzą nie tylko
homomorfizmy, lecz także klasy rozszerzeń.
\end{definicja}

\section{Snop dualizujący i twierdzenie Verdiera}
\begin{definicja}[Snop dualizujący i dualność Verdiera]
\label{def:verdier-v12}
Niech $a:X\to\{*\}$ będzie mapą do punktu.
W formalizmie sześciu operacji dla odpowiedniej kategorii
przestrzeni (np. lokalnie zwartych, lokalnie kontraktowalnych,
o skończonym wymiarze kohomologicznym), $a^!$ oznacza
prawy adjunkt do pochodnego obrazu właściwego $Ra_!$.
Kompleks
$\omega_X^\bullet:=a^!k$ nazywamy \emph{kompleksem dualizującym}.
Dla kompleksu $\mathcal F^\bullet$ definiujemy
\[
\mathbb D_X\mathcal F^\bullet
:=R\mathcal Hom(\mathcal F^\bullet,\omega_X^\bullet).
\]
Przez $D^b_c(X;k)$ oznaczamy kategorię ograniczonych
kompleksów snopów o konstruktybilnych grupach kohomologii,
gdy $X$ jest odpowiednio stratyfikowalna.
\end{definicja}

\begin{twierdzenie}[Globalna dualność Verdiera — wersja nad ciałem]
\label{thm:verdier-global-v12}
Niech $X$ będzie lokalnie zwartą, lokalnie kontraktowalną,
parazwartą przestrzenią o skończonym wymiarze kohomologicznym,
na której istnieje standardowy formalizm $Ra_!\dashv a^!$.
Dla $\mathcal F^\bullet\in D^b_c(X;k)$ zachodzi
naturalny quasi-izomorfizm
\[
R\Gamma\bigl(X,\mathbb D_X\mathcal F^\bullet\bigr)
\simeq R\!\operatorname{Hom}_k
\bigl(R\Gamma_c(X,\mathcal F^\bullet),k\bigr).
\]
W szczególności, gdy $R\Gamma_c(X,\mathcal F^\bullet)$
ma skończeniewymiarowe grupy kohomologii,
$H^{-q}(X,\mathbb D_X\mathcal F^\bullet)
\cong H_c^q(X,\mathcal F^\bullet)^\vee$,
gdzie $V^\vee=\operatorname{Hom}_k(V,k)$.
\end{twierdzenie}
\begin{proof}[Wskazanie dokładnego miejsca użycia adjunkcji]
Z definicji kompleksu dualizującego i z adjunkcji
$Ra_!\dashv a^!$ wynika
$R\!\operatorname{Hom}_X(\mathcal F^\bullet,a^!k)
\simeq R\!\operatorname{Hom}_k(Ra_!\mathcal F^\bullet,k)$.
Lewą stronę identyfikujemy z
$R\Gamma(X,R\mathcal Hom(\mathcal F^\bullet,a^!k))$,
a $Ra_!\mathcal F^\bullet$ z
$R\Gamma_c(X,\mathcal F^\bullet)$.
Po podstawieniu dostajemy zapisaną równoważność.
Dowód istnienia $a^!$ i konstrukcji adjunkcji jest
osobnym, nietrywialnym twierdzeniem teorii pochodnych
snopów, którego tutaj nie wyprowadzamy.
\end{proof}

\begin{twierdzenie}[Postać dualizująca na rozmaitości]
\label{thm:verdier-manifold-v12}
Jeżeli $M$ jest topologiczną $n$-rozmaitością bez brzegu,
orientowalną nad $k$, to
$\omega_M^\bullet\simeq k_M[n]$; bez wyboru orientacji
$\omega_M^\bullet\simeq\operatorname{or}_{M,k}[n]$,
gdzie $\operatorname{or}_{M,k}$ jest lokalnym systemem orientacji.
Dlatego dla $\mathcal F=k_M$ globalna dualność Verdiera
daje dualność Poincarégo z nośnikami zwartymi.
\end{twierdzenie}
\begin{proof}[Wyprowadzenie z lokalnego modelu]
Ustalmy punkt $x\in M$ i kulę współrzędnych $B$.
Lokalna kohomologia z nośnikiem w $x$ spełnia
$H^q(B,B\setminus\{x\};k)\cong k$ dla $q=n$
i znika w innych stopniach. Orientacja lokalna to
wybór generatora tej jednowymiarowej grupy; generatory
w sąsiednich mapach różnią się znakiem stopnia
lokalnej zmiany współrzędnych (albo działaniem $-1$
w ciele $k$). Konstruuje to system $\operatorname{or}_{M,k}$.

Sam rachunek lokalnej kohomologii \emph{nie wystarcza}
do zbudowania globalnego quasi-izomorfizmu. Potrzebujemy
ponadto twierdzenia o lokalizacji dla $a^!$:
po ograniczeniu do $B$ kompleks dualizujący jest
kompleksem dualizującym na $B$, a jego funkcjonalne
odwzorowania przejścia realizują właśnie otrzymany
system orientacji. Z taką identyfikacją lokalną i
zgodnością na przecięciach sklejamy quasi-izomorfizm
$\omega_M^\bullet\simeq\operatorname{or}_{M,k}[n]$.
Twierdzenie o lokalizacji $a^!$ jest tu jawnie
przyjętym składnikiem formalizmu sześciu operacji.
\end{proof}

\begin{przyklad}[Prosta rzeczywista]
Dla $X=\R$ i $\mathcal F=k_\R$ mamy
$\mathbb D_\R(k_\R)\simeq k_\R[1]$.
Stąd $H^{-1}(\R,\mathbb D k_\R)\cong H^0(\R,k)=k$,
co zgadza się z dualną przestrzenią
$H_c^1(\R;k)^\vee\cong k$.
Dla $q=0$ obie strony są zerowe.
\end{przyklad}

\begin{przyklad}[Kiełek a kopiełek snopu stałego na prostej]
\label{prz:verdier-real-line-costalk-23}
Niech $i:\{0\}\hookrightarrow\R$ będzie włożeniem punktu,
a $k_\R$ snopem funkcji lokalnie stałych o wartościach w ciele
$k$. Zwykły kiełek $i^*k_\R$ to $k$ umieszczone w stopniu
$0$. Nadzwyczajny obraz odwrotny jest natomiast obliczany
przez kohomologię lokalną:
\[
 H^j(i^!k_\R)\cong H^j(I,I\setminus\{0\};k),
 \qquad I=(-1,1).
\]
Ponieważ $I$ jest ściągalny, a $I\setminus\{0\}$ ma
dwie składowe spójności, początek ciągu względnego ma postać
\[
 0\longrightarrow H^0(I,I\setminus\{0\};k)
 \longrightarrow k\xrightarrow{a\mapsto(a,a)} k^2
 \longrightarrow H^1(I,I\setminus\{0\};k)
 \longrightarrow0.
\]
Przekątna jest różnowartościowa, a jej kokernelem jest $k$,
więc $i^!k_\R\simeq k[-1]$ w kategorii pochodnej,
podczas gdy $i^*k_\R\simeq k$.
Jest to konkretny przypadek różnicy między operacjami
$i^*$ i $i^!$: lokalne dane o nośnikach są przesunięte
o rzeczywisty wymiar rozmaitości.
Zgodnie z Twierdzeniem~\ref{thm:verdier-manifold-v12}
kompleks dualizujący na prostej jest równy $k_\R[1]$.
W szczególności lokalny rachunek stopni zgadza się
z przesunięciem w formule dualności Verdiera.
\end{przyklad}

\begin{przyklad}[Węzeł krzywej zespolonej]
\label{ex:node-duality-v12}
Niech $X=\{(z,w)\in\C^2:zw=0\}\cap B_\varepsilon$,
gdzie $B_\varepsilon$ jest małą otwartą kulą wokół $0$.
Otoczenie węzła jest topologicznie stożkiem na
dwóch rozłącznych okręgach (geometryczny link w $S^3$
jest splotem Hopfa), a $X\setminus\{0\}$ deformacyjnie zwija się
do $S^1\sqcup S^1$. Ponieważ $X$ jest kontraktowalna,
ciąg kohomologii pary daje
\[
H^q(X,X\setminus\{0\};k)\cong
\widetilde H^{q-1}(S^1\sqcup S^1;k)
=\begin{cases}
k,&q=1,\\k^2,&q=2,\\0,&q\notin\{1,2\}.
\end{cases}
\]
Natomiast w punkcie gładkim zespolonej krzywej
lokalna kohomologia pary jest niezerowa tylko w stopniu $2$,
gdzie wynosi $k$. Dlatego \emph{stały kompleks przesunięty}
$k_X[2]$ nie może odgrywać roli kompleksu dualizującego
w takiej samej postaci jak na rozmaitości gładkiej.
To pierwszy sygnał pojawienia się w geometrii
algebraicznej kompleksów dualizujących i snopów perwersyjnych.
\end{przyklad}

\section{Sześć operacji i właściwe odwzorowania}
Na poziomie pochodnym porządkujemy operacje
$f^{-1}$ (odwrotny obraz snopu), $Rf_*$ (obraz bezpośredni),
$Rf_!$ (obraz właściwy), $f^!$ (nadzwyczajny odwrotny obraz),
$\otimes^L$ i $R\mathcal Hom$.
Dla snopów \emph{wektorowych nad ciałem} zwykłe $f^{-1}$
i $\otimes$ są już dokładne, dlatego w tej kategorii $f^{-1}$ nie wymaga
pochodnego poprawiania, a $\otimes^L=\otimes$.
Zwykłą konstrukcję iloczynu tensorowego i jej własność uniwersalną
poznaliśmy w
\href{https://github.com/dariusz-piekarz/mathematical_scipts/blob/biblioteka-matematyczna-iteracja-6/chapters/pl/differential_geometry_and_topology/03-algebra-tensorowa-i-zewnetrzna.tex}{rozdziale 3 skryptu \emph{Geometria różniczkowa}:
\emph{Algebra tensorowa i zewnętrzna}}.
% Historyczne etykiety DG: ch:algebra-tensorowa,
% sec:iloczyn-tensorowy, thm:tensor-hom.
W obecnym rozdziale używamy tej konstrukcji dla
snopów $k$-przestrzeni wektorowych w poszczególnych
otwartych zbiorach, a następnie przechodzimy do kompleksów.
Równość z \emph{pochodnym} iloczynem tensorowym
wynika tu z płaskości przestrzeni wektorowych nad ciałem;
nie jest ogólną własnością modułów nad dowolnym pierścieniem.
Nie powtarzamy wcześniejszego dowodu własności uniwersalnej.

\begin{stwierdzenie}[Obraz właściwy i dualność]
\label{prop:verdier-proper-v12}
Jeśli $f:X\to Y$ jest właściwym ciągłym odwzorowaniem
między przestrzeniami spełniającymi hipotezy
Twierdzenia~\ref{thm:verdier-global-v12}, to
$Rf_!\cong Rf_*$. Dla konstruktybilnego
$\mathcal F^\bullet$ zachodzi ponadto naturalna
równoważność
\[
Rf_*\,\mathbb D_X\mathcal F^\bullet
\simeq\mathbb D_Y\,Rf_*\mathcal F^\bullet.
\]
\end{stwierdzenie}
\begin{proof}[Argument z formalizmu sześciu operacji]
Właściwość mapy oznacza, że warunek właściwości
nośnika sekcji na przeciwobrazie dowolnego otwartego
zbioru jest automatyczny; więc $f_!=f_*$ i tak samo
po przejściu do funktorów pochodnych. Zgodność z dualnością
otrzymujemy z twierdzenia o dualności względnej
$Rf_*R\mathcal Hom_X(\mathcal F,f^!\mathcal G)
\simeq R\mathcal Hom_Y(Rf_*\mathcal F,\mathcal G)$
razem z $(a_Y\circ f)^!\simeq f^!a_Y^!$.
Jest to zastosowanie twierdzenia o dualności względnej,
a nie jego niezależny dowód.
\end{proof}

\begin{uwaga}[Co wchodzi do geometrii algebraicznej?]
Dla przestrzeni analitycznej $X^{\mathrm{an}}$ mamy snop
stały $k_X$ i jego kohomologię topologiczną. Mamy też
snop funkcji holomorficznych $\mathcal O_X$ oraz dla
schematów snopy koherentne w topologii Zariskiego.
Nie należy bez hipotez identyfikować
$H^*(X^{\mathrm{an}},k)$, $H^*(X,\mathcal O_X)$
i kohomologii etalnej. Należą do różnych teorii
współczynników i, w tym ostatnim przypadku,
topologii Grothendiecka. Dualność Verdiera prowadzi
od lokalnych systemów do kompleksów konstruktybilnych,
a dalej do kompleksów przecięcia i twierdzenia o rozkładzie.
Tych zaawansowanych wyników nie uznajemy jeszcze za udowodnione.
\end{uwaga}

\section{Ciąg lokalizacji dla podzbioru zamkniętego}
\begin{twierdzenie}[Lokalizacja nośników]
\label{thm:localization-v12}
Niech $X$ będzie lokalnie zwartą przestrzenią Hausdorffa,
$Z\subseteq X$ zbiorem domkniętym, $U=X\setminus Z$,
a $j:U\hookrightarrow X$, $i:Z\hookrightarrow X$
kanonicznymi włożeniami. Dla snopu $\mathcal F$
$k$-przestrzeni wektorowych na $X$ istnieje naturalny
ciąg dokładny
\[
\cdots\to H_c^q(U,\mathcal F|_U)\to H_c^q(X,\mathcal F)
\to H_c^q(Z,\mathcal F|_Z)
\to H_c^{q+1}(U,\mathcal F|_U)\to\cdots.
\]
\end{twierdzenie}
\begin{proof}
Rozpatrzmy ciąg snopów na $X$:
\[
0\longrightarrow j_!(\mathcal F|_U)
\longrightarrow\mathcal F\longrightarrow i_*(\mathcal F|_Z)
\longrightarrow0.
\]
Sprawdźmy dokładność, nie ograniczając się do sekcji globalnych.
Dla $x\in U$ na włóknach otrzymujemy
$0\to\mathcal F_x\xrightarrow{\id}\mathcal F_x\to0$.
Dla $x\in Z$ otrzymujemy
$0\to0\to\mathcal F_x\xrightarrow{\id}\mathcal F_x\to0$,
ponieważ ograniczenie snopu do $Z$ nie zmienia
kiełków przy punktach $Z$. Zatem ciąg jest dokładny
na mocy Lematu~\ref{lem:stalk-exact-v12}.

Pochodny funktor sekcji o nośnikach zwartych
przekształca ten ciąg w długi ciąg kohomologii.
Pierwszy człon identyfikujemy ze wspomnianą
kohomologią $U$ przez dokładność $j_!$.
Ponieważ $i$ jest zanurzeniem domkniętym i właściwym,
$\Gamma_c(X,i_*\mathcal G)=\Gamma_c(Z,\mathcal G)$;
identyfikacja pozostaje prawdziwa po zastosowaniu
funktorów pochodnych. Otrzymujemy wzór z tezy.
\end{proof}

\begin{przyklad}[Usunięcie punktu z prostej]
Dla $X=\R$, $Z=\{0\}$ i $U=\R\setminus\{0\}$ mamy
$H_c^0(Z;k)=k$, $H_c^1(U;k)=k^2$ oraz
$H_c^1(\R;k)=k$. Istotny fragment ciągu lokalizacji to
\[
0\longrightarrow k\xrightarrow{\Delta}k^2
\xrightarrow{\,(a,b)\mapsto b-a\,}k\longrightarrow0,
\]
przy odpowiednim wyborze generatorów orientacji
obu półprostych. Łączący homomorfizm $\Delta$
jest przekątny: usunięcie jednego punktu tworzy dwie
niezależne klasy o nośnikach zwartych w stopniu pierwszym.
\end{przyklad}

\section{Przykład z geometrii algebraicznej: kohomologia \texorpdfstring{$\mathcal O(m)$}{O(m)} na \texorpdfstring{$\mathbb P^1$}{P1}}
\label{sec:p1-sheaf-v12}
Nie poprzestawajmy na formalizmie: snopy pozwalają
wykonać rachunek, którego nie opisuje zwykła homologia
singularna. Niech $\mathbb P^1=\mathbb P^1(\C)$,
$U_0=\{[x:y]:x\ne0\}$, $U_\infty=\{[x:y]:y\ne0\}$.
Oznaczmy $z=y/x$ na $U_0$, $w=x/y$ na $U_\infty$.
Na przecięciu $w=1/z$.

\begin{definicja}[Snop $\mathcal O(m)$]
Dla $m\in\Z$ wiązkę liniową $\mathcal O(m)$
trywializujemy na $U_0$ i $U_\infty$ przez
lokalne ramy $e_0$ i $e_\infty$ z regułą przejścia
$e_\infty=z^m e_0$ na $U_0\cap U_\infty$.
Snop $\mathcal O(m)$ oznacza snop \emph{holomorficznych}
przekrojów tej wiązki.
\end{definicja}

\begin{twierdzenie}[Kohomologia skręceń na prostej rzutowej]
\label{thm:p1-twist-v12}
Dla każdego $m\in\Z$ zachodzą izomorfizmy
\[
H^0(\mathbb P^1,\mathcal O(m))\cong
\begin{cases}\C^{m+1},&m\ge0,\\0,&m<0,\end{cases}
\qquad
H^1(\mathbb P^1,\mathcal O(m))\cong
\begin{cases}\C^{-m-1},&m\le-2,\\0,&m\ge-1.\end{cases}
\]
Ponadto $H^q(\mathbb P^1,\mathcal O(m))=0$ dla $q\ge2$.
\end{twierdzenie}
\begin{proof}
Zaczniemy od kompleksu \v{C}echa dla pokrycia
$\{U_0,U_\infty\}$, identyfikując przekroje z funkcjami
holomorficznymi w odpowiednich trywializacjach.
Zgodność przekrojów na przecięciu ma postać
$f_0(z)=z^m f_\infty(1/z)$.
Zatem kompleks ma dwa niezerowe wyrazy:
\[
0\longrightarrow
\mathcal O(\C)\oplus\mathcal O(\C)
\xrightarrow{\,\delta\,}\mathcal O(\C^*)
\longrightarrow0,
\qquad
\delta(f_0,f_\infty)=f_0(z)-z^m f_\infty(1/z).
\]
Jego jądro liczy globalne przekroje: funkcja
$f_0\in\mathcal O(\C)$ może przedłużyć się zgodnie
z ramą w nieskończoności wtedy i tylko wtedy,
gdy jest wielomianem stopnia co najwyżej $m$ przy $m\ge0$;
przy $m<0$ pozostaje jedynie przekrój zerowy.

Obliczmy teraz kokernele. Każda funkcja holomorficzna
na $\C^*$ ma rozwinięcie Laurenta
$h(z)=\sum_{j\in\Z}a_jz^j$, zbieżne lokalnie jednostajnie
na $\C^*$. Część $\sum_{j\ge0}a_jz^j$ jest funkcją
całkowitą na $\C$ i może zostać usunięta przez
odjęcie $f_0$. Część $\sum_{j\le m}a_jz^j$
ma postać $z^m f_\infty(1/z)$, gdzie
$f_\infty(w)=\sum_{j\le m}a_jw^{m-j}$ jest
funkcją całkowitą w $w$; można więc usunąć również
wszystkie potęgi $j\le m$. Gdy $m\ge-1$, te dwa
zbiory wykładników pokrywają całą $\Z$.
Gdy $m\le-2$, pozostają dokładnie klasy
$z^{m+1},z^{m+2},\ldots,z^{-1}$, których liczba wynosi
$-m-1$. Stąd obliczenie $H^1$ kompleksu \v{C}echa.

Aby utożsamić je z \emph{pochodną} kohomologią snopu,
stosujemy Twierdzenie~\ref{thm:leray-sheaf-v12}:
$U_0\cong\C$, $U_\infty\cong\C$ i
$U_0\cap U_\infty\cong\C^*$ są przestrzeniami Steina,
a kohomologie wyższych stopni snopów holomorficznych
przekrojów wiązek wektorowych znikają na nich
na mocy \emph{twierdzenia B Cartana}. To ostatnie
jest tutaj wyraźnie przywołanym wynikiem analizy
zespolonej, a nie ukrytą konsekwencją samej definicji
snopu. Pokrycie jest zatem acykliczne i rachunek
\v{C}echa daje wszystkie grupy kohomologii.
\end{proof}

\begin{uwaga}[Dwa różne rodzaje kohomologii]
Na $\mathbb P^1$ singularna kohomologia ze stałymi
współczynnikami ma $H^0\cong H^2\cong\Z$ i
$H^1=0$, niezależnie od $m$. Natomiast
$H^1(\mathbb P^1,\mathcal O(-3))\cong\C^2$.
To nie jest sprzeczność: grupy badają ten sam
zbiór topologiczny, ale \emph{inne współczynniki snopowe}.
\end{uwaga}


\begin{przyklad}[Parowanie Serre'a na $\mathbb P^1$ obliczone przez residuum]
\label{prz:p1-serre-residue}
Niech $m\le-2$, a $N=-m-2\ge0$.
W dowodzie Twierdzenia~\ref{thm:p1-twist-v12}
otrzymaliśmy bazę
\[
 H^1(\mathbb P^1,\mathcal O(m))
 =\operatorname{span}_{\C}
   \{[z^{m+1}],[z^{m+2}],\ldots,[z^{-1}]\},
\]
podczas gdy
$H^0(\mathbb P^1,\mathcal O(N))$
jest przestrzenią wielomianów stopnia co najwyżej $N$.
Dla klasy $[h]\in H^1(\mathbb P^1,\mathcal O(m))$
i wielomianu $p(z)=\sum_{\ell=0}^N b_\ell z^\ell$
zdefiniujmy
\[
 \langle[h],p\rangle
  =\operatorname{Res}_{z=0}
     \bigl(h(z)p(z)\,\dd z\bigr),
\]
czyli współczynnik przy $z^{-1}$ w szeregu Laurenta
iloczynu $hp$. Jest to dobrze określone
\emph{na klasach kohomologii}.
Istotnie, zmiana reprezentanta Czecha
$h\mapsto h+f_0(z)-z^m f_\infty(1/z)$
z $f_0,f_\infty\in\mathcal O(\C)$ nie zmienia residuum:
$p f_0$ zawiera wyłącznie nieujemne potęgi $z$,
a $z^m p(z)f_\infty(1/z)$ zawiera wyłącznie
potęgi co najwyżej $m+N=-2$.
Żaden z tych składników nie zawiera $z^{-1}$.

Dla baz $[z^{m+1}],\ldots,[z^{-1}]$
i $1,z,\ldots,z^N$ parowanie ma macierz
z jedynkami na przeciwnej przekątnej:
\[
 \bigl\langle[z^{m+1+i}],z^\ell\bigr\rangle
 =\begin{cases}
  1,&\ell=N-i,\\
  0,&\ell\ne N-i,
 \end{cases}
 \quad 0\le i,\ell\le N.
\]
Jest więc niedegenerujące i daje izomorfizm
\[
 H^1(\mathbb P^1,\mathcal O(m))
 \cong H^0(\mathbb P^1,\mathcal O(-m-2))^\vee.
\]
Na przykład przy $m=-3$ bazami są
$([z^{-2}],[z^{-1}])$ i $(1,z)$,
a macierz parowania wynosi
$\left(\begin{smallmatrix}0&1\\1&0\end{smallmatrix}\right)$.
Dla $m=-1$ obie strony są zerowe;
dla $m\ge0$ także znikają, bo
$H^1(\mathcal O(m))=0$
oraz $H^0(\mathcal O(-m-2))=0$.
Zatem wzór pozostaje prawdziwy dla każdego $m\in\Z$.

Przeprowadziliśmy w ten sposób \emph{konkretny rachunek}
przypadku dualności Serre'a na prostej rzutowej,
ale nie ogólny dowód dualności Serre'a
dla dowolnej zwartej rozmaitości zespolonej.
\end{przyklad}

\begin{figure}[!htbp]
\centering
\begin{tikzpicture}[>=Stealth,font=\small]
 \node[draw,rounded corners=4pt,fill=manifoldblue!8,
       minimum width=3.9cm,minimum height=1.1cm,
       align=center] (a) at (0,0)
       {$H^1(\mathbb P^1,\mathcal O(-3))$\\
        $[z^{-2}],\,[z^{-1}]$};
 \node[draw,rounded corners=4pt,fill=manifoldgreen!8,
       minimum width=3.9cm,minimum height=1.1cm,
       align=center] (b) at (6.0,0)
       {$H^0(\mathbb P^1,\mathcal O(1))$\\
        $1,\,z$};
 \draw[->,thick] (a)--(b)
       node[midway,above=7pt] {residuum};
 \node[draw,rounded corners=4pt,fill=manifoldgold!12,
       minimum width=4.8cm,minimum height=1.05cm,
       align=center] (c) at (3,-2)
       {$\langle[z^{-2}],z\rangle=1$\\
        $\langle[z^{-1}],1\rangle=1$};
 \draw[->,thick] (a.south east)--(c.north west);
 \draw[->,thick] (b.south west)--(c.north east);
\end{tikzpicture}
\caption{Parowanie na $\mathbb P^1$ obliczone
przez współczynnik przy $z^{-1}$.
W wybranych bazach jego macierz ma jedynki
na przeciwnej przekątnej.}
\label{fig:p1-serre-residue}
\end{figure}

\section{Podsumowanie, ograniczenia i dalszy kierunek}
W tym rozdziale udowodniliśmy bez aparatu pochodnego
warunek snopowy, sprawdzanie dokładności na włóknach,
obliczenie $H^0$ metodą \v{C}echa, opis przedłużenia przez zero
i elementarne przykłady. Dla kohomologii z nośnikami
wyprowadziliśmy ciąg Mayera--Vietorisa z dokładnego
ciągu snopów. Rezolwenta de Rhama tłumaczy,
dlaczego lokalny lemat Poincarégo daje globalne
kohomologie z nośnikami zwartymi.

Dualność Verdiera jest tu sformułowana precyzyjnie
\emph{względem} istnienia funktora $a^!$
i standardowego formalizmu sześciu operacji.
Nie rozwinęliśmy niezależnego dowodu jego istnienia,
pełnej teorii konstruktybilności ani
pełnego dowodu względnej dualności. Takie rozdzielenie
statusu dowodów będzie ważne przy scalaniu skryptu
z teorią snopów i geometrią algebraiczną.


